% ============================================================
%  H. Høffding, "Forholdet mellem Tro og Viden i dets
%  historiske Udvikling. Et indledende Foredrag."
%  Først trykt 1885; optrykt i: Mindre Arbejder,
%  Kjøbenhavn: Gyldendal, 1899.
%  TRANSCRIPTION (Danish)
% ============================================================
\documentclass[12pt,a4paper]{article}
\usepackage[utf8]{inputenc}
\usepackage[T1]{fontenc}
\usepackage[danish]{babel}
\usepackage{microtype}
\usepackage{setspace}
\usepackage[margin=1.2in]{geometry}
\usepackage{fancyhdr}
\usepackage{hyperref}

\hypersetup{
  pdftitle={Forholdet mellem Tro og Viden i dets historiske Udvikling},
  pdfauthor={H. Høffding},
  hidelinks
}

\pagestyle{fancy}
\fancyhf{}
\rhead{Høffding, \emph{Forholdet mellem Tro og Viden} (1885)}
\cfoot{\thepage}

\setstretch{1.3}

\title{\textbf{Forholdet mellem Tro og Viden\\[4pt]
       i dets historiske Udvikling}\\[6pt]
  \large Et indledende Foredrag}
\author{Harald Høffding}
\date{Trykt i: \emph{Mindre Arbejder}, Kjøbenhavn: Gyldendal, 1899\\
  (Foredraget holdt 1885)}

\begin{document}
\maketitle
\thispagestyle{fancy}

\bigskip

\noindent Problemet om Tro og Viden opstaar med Nødvendighed paa et vist
Trin af Kulturudviklingen. Vi træffe det saaledes inden for den indiske,
den græske og den muhammedanske Verden, saa snart den aandelige, særligt
den intellektuelle Udvikling har naaet et vist Punkt. Her ville vi dog
indskrænke os til at fremhæve de Hovedformer, under hvilke det er
fremtraadt inden for den evropæiske, ved Kristendommen paavirkede
Aandsudvikling.

\section*{I.}

Af Kristendommens Kildeskrifter kan ingen Belæring hentes angaaende dette
Spørgsmaal. Der eksisterer paa det Standpunkt, ud fra hvilket de ere
forfattede, ingen fra Troen forskellig Viden. Blikket er rettet mod én
eneste Genstand, et eneste Maal, der lægger Beslag paa al Stræben og al
aandelig Kraft. Bestræbelsen for at erkende Sandheden er her umiddelbart
Et med Trangen til at vinde Sjælens Frelse og gaar ikke sin egen Vej med
sin egen Metode. Der er ingen Forskel mellem Teori og Praksis, da Teori og
teoretisk Interesse aldeles ikke eksisterer. Naar Guddommen selv har
aabenbaaret sig synligt for Alle, --- naar det ved Troen er muligt at naa
til Forstaaelse af det Mysterium, som indeslutter alle Visdommens og
Kundskabens Skatte, --- saa mangler ethvert Motiv til at søge Kundskab ad
andre Veje.

Det var naturligt, at Troen fra først af maatte være Et og Alt, saa der
ikke blev Plads til andre aandelige Fornødenheder. Den første Slægt stod
endnu saa nær ved den mægtige aandelige Bevægelse, der førte til den
kristelige Menigheds Grundlæggelse, at den fuldstændigt levede og aandede
i de store Minder og Syner. Dertil kom endnu, at denne første Slægt levede
i Troen paa Kristi snarlige Genkomst. Da skulde Alt afsluttes,
Verdensdommen fældes, og alle menneskelige og jordiske Forhold give Plads
for en højere Tingenes Orden. Hvem havde med en saadan Katastrofe for Øje
Tid og Ro i Sindet til filosofisk Tænkning eller videnskabelig Forsken?
Videnskaben kunde lige saa lidt som nogen anden menneskelig Bestræbelse faa
nogen Betydning. Der skulde jo slet ikke finde nogen historisk Udvikling
Sted; Maalet var naat, og Enden stod for Døren. Man forstaar ikke det ny
Testamente og dets Forhold til Kulturbestræbelsen i dens forskellige
Former, naar man ikke fastholder dette. Det gik her med den videnskabelige
Stræben som med Stræben efter personlig Frihed og med Indgaaelse af
Ægteskab (se 1.\ Korinterbrev 7.\ Kap.). Hvorledes kunde alt Sligt faa
Betydning, eller endog blot blive Genstand for Drift, naar alle Tanker og
Kræfter maatte anspændes for at være forberedt til den store, den absolute
Afgørelse!

\section*{II.}

Tiden gik, Afslutningen kom ikke. Kirken trøstede sig med, at for Gud er
tusind Aar som én Dag, og én Dag som tusind Aar (2.\ Peters Brev 3.\ Kap),
og begyndte snart at falde til Ro i Verden, at indrette sig efter
Forholdene og at tage Stilling over for de forskellige Sider af de
kulturhistoriske Bestræbelser. Over for Videnskaben blev dens Stilling den,
at den ikke længer i Troen fandt al fornøden Erkendelse, men søgte at vise,
at Troens Lærdomme satte Kronen paa al menneskelig Erkendelse. Kirken
tilegnede sig den Videnskab, Grækerne havde udviklet: den græske Astronomi,
Medicin og Filosofi, og anbragte sin Teologi som Spiret paa Videnskabens
Bygning, --- paa lignende Maade som den omdannede de græske Templer til
Kirker og f.\ Eks.\ indviede Partenon, Athenes Tempel, til den »hellige
Visdom«.

Allerede i det saakaldte Johannesevangelium finde vi efter Nogles Mening
platonisk Filosofi og kristelig Tro forbundne. Denne Bestræbelse fortsættes
af Kirkefædrene og naar sin Afslutning i Teologiens klassiske Periode hos
den middelalderlige Kirkes store Lærere. Hos Thomas fra Aqvino (født 1227,
død 1274), den katolske Kirkes endnu i vore Dage hædrede og anerkendte
Lærer, stiller Forholdet mellem Tro og Videnskab sig paa følgende Maade.
Videnskaben udforsker alle underordnede og specielle Aarsager, medens den
Troendes Tanke angaar den første, den absolute Aarsag til Alt. Kun
Teologien lærer os derfor ret at forstaa Tingene; Videnskaben lærer os
egentligt kun, hvad Tingene ere, ikke hvorfra de stamme. Men der er ingen
Strid mellem Troen, der er Teologiens Princip, og Fornuften, der er
Videnskabens Princip. Fornuften stammer fra Gud lige saa vel som Troen, og
Troen strider ikke mod Fornuften, skønt den gaar ud over Fornuften.

I poetisk Form finde vi dette karakteristiske middelalderlige Standpunkt
hos Dante (født 1265, død 1321), som var en Discipel af Thomas. I sin
»Comedia« beretter Digteren, hvorledes det, for at hans Sjæl kunde blive
lutret, forundes ham at gennemvandre den hinsidige Verden: Helvede,
Skærsilden og Paradiset. Det er hans forklarede Elskede, Beatrice, Symbolet
for den paa himmelsk Aabenbaring grundede Teologi, som skænker ham denne
Gunst og sender Vergil, Symbolet for den naturlige Videnskab, til ham for
at lede ham gennem Underverdenen. Men Himlen tør Vergil ikke betræde; han
angiver selv, med stor Sagtmodighed, Grunden:

\begin{quote}
ti Herskeren, hvem Himlens Stad tilhører,\\
vil ej, at jeg dens Tærskel maa betræde,\\
fordi jeg mod hans Love var Oprører.
\end{quote}

Da de Vandrende derfor nærme sig Himlen, aabenbarer Beatrice sig selv, og
Vergil forsvinder.

Saa from og ydmyg var den naturlige Fornuft bleven i Middelalderen. Men
dens Kræfter voksede, og Middelalderens Slutning betegnes ved voldsom Splid
mellem Beatrice og Vergil, Tro og Viden. Vi træffe her Problemet stillet
for første Gang.

\section*{III.}

Den stolte Bygning, som Thomas fra Aqvino havde fuldført, og som Dante
havde besunget, vaklede snart i sin Grundvold. Den fortsatte Grublen
opdagede, at den var bygget af Elementer, der ikke kunde bringes i Harmoni.
Allerede Vilhelm af Occam (\dag\ 1347) finder en Modstrid mellem den
naturlige Fornuft og den over enhver Lov hævede guddommelige Almagt. Han
lærer, at Noget kan være sandt i Teologien, som ikke er det i Filosofien:
dersom Bibelen og Kirken lære Noget, som er selvmodsigende, maa det
alligevel tages for Sandhed! Teologien bygger paa overnaturlig Avtoritet og
er derfor ingen Videnskab.

Have vi da to Slags Erkendelse, to Slags Logik? Dette Spørgsmaal, der
ligger saa nær, bliver ikke fremdraget. Occam selv var troende Katolik; han
paaviste Spliden, men forfulgte ikke dens Konsekvenser. Den katolske Kirke
holdt fast ved sin hellige Thomas og har ikke tilladt nogen anden Opfattelse
af Tro og Viden at gøre sig gældende inden for dens Enemærker. Endnu i vort
Aarhundrede ere ikke mindre end tre katolske Filosofer blevne dømte, fordi
de enten vare mere eller mindre rationalistiske end St.\ Thomas. For den
katolske Kirke eksisterer der her ikke længer noget Problem. --- Lad os da
se, hvorledes Sagen stiller sig inden for Protestantismen.

\section*{IV.}

I den nyere Tid breder Livet sig i saare mange forskellige Retninger og har
et langt mere broget og mangesidigt Præg end i Oldtid og Middelalder. Og
Kirken staar kun som ét af de mange Omraader for Livsbevægelsen; den er
ikke længer Et og Alt som i den første Menighed, eller Kronen paa det Hele
som i Middelalderen. Historien er ikke først og fremmest en Kirkehistorie,
men blandt Andet ogsaa Kirkehistorie. Overvejende passivt modtager Kirken
de nye Livsretningers Tilbageslag og søger at finde sig til Rette over for
dem; men de udgaa ikke længer fra den.

Den kirkelige Protestantismes Forhold til Videnskaben er omtrent det samme
som Katolicismens. Luther havde vel beraabt sig --- ej blot paa den hellige
Skrift, men ogsaa paa »klare Grunde«; men ellers anvendte han paa Fornuften
Pauli Ord, at Kvinden skal tie i Menigheden. Altsaa Fornuften skulde have
Vergils Rolle, maatte ikke betræde det Allerhelligste. Den protestantiske
Kirke vilde gerne vedblive at spille »den guddommelige Komedie«. Men Vergil
har ikke længer vist sig medgørlig.

Den protestantiske Ortodoksi blev en ny Skolastik, men mere ængstelig og
mindre storartet end den middelalderlige. --- Den har ikke fundet nogen
Dante. De store protestantiske Digtere gaa uden om den eller imod den.
Poesiens Forhold til Kirken danner en Parallel til Videnskabens: Oldkirken
har ingen Poesi, Katolicismen en kirkelig Poesi, Protestantismen en uden for
Kirken staaende, om ikke antikirkelig Poesi.

For Protestantismen er det 18.\ Aarhundrede, hvad det 14.\ var for
Katolicismen. Men den Krise, som det bragte, var langt grundigere, bundede
langt dybere end Renaissancebevægelsen. Det er den selvstændige humane
Livsanskuelses Fødselstid. Forløberne vare naturligvis mange --- men først
nu forstodes og anerkendtes de. Først nu fandt f.\ Eks.\ Spinoza, der i
halvandet Aarhundrede havde ligget »som en død Hund«, en Kreds, der kunde
tænke hans Tanker og tilegne sig dem. Først nu fik den teologiske
Livsanskuelse en jævnbyrdig Modstander. Jeg tænker her ikke paa noget
bestemt, afsluttet System, men paa en almindelig Tanke- og Livsretning, som
i de store Træk var fælles for Mænd som Lessing og Kant, Schiller og Goethe.
--- Som man ser, fandt denne nye Livsanskuelse straks mere end én Dante.

Det Ejendommelige ved dette Standpunkt, hvis Hovedtanker kun forbigaaende
ere blevne stillede i Skygge ved den romantiske Reaktion, er --- fra den
teoretiske Side set ---, at man over for den teologiske Tro ikke stiller en
foregiven Indsigt i de teologiske Lærdommes Umulighed, men en Paastand om
deres Ubevislighed. Kants store Gerning var, fra denne Side set, at give en
videnskabelig Lære om Videnskabens Grænser. Ikke ydmyget og kuet, som en
Oprører, der er falden til Føje, ikke ude fra hæmmet, men med selvstændig
og klar Selverkendelse siger den menneskelige Fornuft her: saa langt kan jeg
naa, og ikke videre! --- Erkendelsens Grænser følge af Erkendelsens Natur.
Ved Undersøgelse af vor Erkendelses sidste Forudsætninger og Love afstikkes
det Omraade, paa hvilket den kan bevæge sig. Dette Omraade er uoverskueligt,
men dog bestemt afgrænset over for det, der efter sin Natur ikke kan blive
Genstand for vor Erkendelse, saaledes som denne nu én Gang er beskaffen.

For Kant og hans Tid laa Erkendelsens Betydning ikke i, at den skulde give
os »den absolute Sandhed«, men deri, at vi ved den stedse kunne gaa frem i
Sandhed. Den evige Søgen, den evige Stræben blev fra nu af Slagordet:

\begin{quote}
\itshape
Wer immer strebend sich bemüht,\\
den können wir erlösen.
\end{quote}

Ikke som om Intet kunde findes eller opnaas: men ethvert Fund, ethvert
opnaaet Maal staar som Udgangspunkt for ny Stræben. Det var en stor Livs-
og Verdenslov, man her gav Udtryk, --- Udviklingstanken i en af dens
vigtigste Former, en Tanke, som vor Tid har givet en realistisk Underbygning
og en nærmere kritisk Bestemmelse, men som det er hin Slægts Ære at have
udtalt for første Gang med fuld og klar Bevidsthed.

Men endnu vigtigere var et andet Princip, der blev slaaet fast som
Hovedpunkt i den humane Livsanskuelse: den frie Samvittigheds Princip.
Kants anden store Gerning var, at han hævdede Muligheden og Nødvendigheden
af en rent human Etik, fastslog, at Samvittigheden lige saa lidt som den
teoretiske Forsken kan lade sig sine Grænser afstikke ude fra. Maalestokken
for Alt, hvad vi skulle anerkende som godt, ædelt og ophøjet, maa findes i
vort eget Indre og være bestemt ved den menneskelige Natur.

Naar de, der staa paa dette Standpunkt, ikke kunne gaa ind paa den kirkelige
Tro, saa er det, fordi den strider ej blot mod deres Forstand, men mod deres
Samvittighed. Man skal ikke blot slutte sig til en én Gang for alle given
Formulering af Sandheden, men ogsaa bøje sig for Lærdomme, der ikke kunne
bringes til at stemme med Samvittighedens Maalestok for Godt og Ondt.

Først her bliver Problemet skærpet, og dets egentlige Brod træder for
Dagen. Hvad der fra den ene Side gøres til den højeste, paa en Maade den
eneste Pligt, en Pligt, hvis Ikke-Opfyldelse medfører Fortabelse, --- det
afvises fra den anden Side netop som stridende mod Samvittighedens klare og
bestemte Fordring! ---

Biskop Martensen har fra sit Standpunkt klart set, at det vigtigste
Vendepunkt i Striden om Tro og Viden fremkom ved den etiske Humanismes
Grundlæggelse i Slutningen af forrige Aarhundrede. Han beskæftiger sig i sin
»Kristelige Etik« vidtløftigt med denne Retnings Repræsentanter og deres
Forhold til Kirkens Tro. Den sidste Forklaring af deres Vægring ved at
slutte sig til denne finder han med Rette ikke i deres Forstand, men i deres
Villie --- kun at denne Villie efter hans Paastand er Trods mod det Hellige:
»Mod Guds Hellighed have de en hemmelig Antipati!«

Hvor underligt, eller rettere hvor forfærdeligt maa det menneskelige Liv dog
ikke tage sig ud fra det ortodokse Standpunkt! Den redeligste Søgen og
Stræben, den strengeste Samvittighedsfuldhed forenet med den mest
uforbeholdne Anerkendelse af Ufuldkommenheden i det, der er naat, maa fra
dette Standpunkt udlægges som Antipati mod Helligheden. --- Har man ikke
forstaaet, hvorfor Problemet om Tro og Viden saa let sætter Lidenskaberne i
Bevægelse, saa forstaar man det her.

\section*{V.}

Det nye humane Standpunkt syntes at have ret god Udsigt til at trænge
igennem. Tidsalderens mest fremragende Aander hyldede det, de ledende
Statsmænd virkede i den af det betegnede Retning, i selve Kirken var
Ortodoksien afgjort svækket, og de konfessionelle Forskelligheder traadte
tilbage for en oplyst og frisindet Tankegang, der var fælles for Alle.
Saaledes vare ved Aarhundredets Begyndelse Forholdene ikke blot i den
protestantiske, men ogsaa i den katolske Kirke.

Men hvor langt anderledes skulde Forholdene ikke udvikle sig i Løbet af
Aarhundredet, og hvor forandrede vare de ikke allerede efter de to første
Tiaars Udløb!

Et enkelt Træk viser tydeligt Forandringen. I Aaret 1799 udgav Schleiermacher
sine »Taler om Religionen til de Dannede blandt dens Foragtere«. Han kæmpede
her ingenlunde for nogen dogmatisk Tro, men hævdede en panteistisk
Religiøsitet og forherligede Spinoza. For en saadan Anskuelse vilde han vinde
Anerkendelse hos de Dannede. Men da han i Aaret 1821 besørgede den tredje
Udgave af sin Bog, erklærede han i Fortalen, at naar man nu saa sig om blandt
de Dannede, vilde man snarere finde det fornødent at skrive Taler til Heldører
og Bogstavtrælle, uvidende og ukærlige Fordemmere, Overtroiske og
Hyperdogmatikere!

Fire Aar i Forvejen (1817) havde den holstenske Præst Claus Harms til Minde
om Luthers Optræden udgivet 95 nye Teser, i hvilke han blandt Andet
protesterer mod »vor Tids Paver«: Fornuften paa det teoretiske Omraade, og
Samvittigheden paa det praktiske Omraade. --- Hvor var det dog muligt, at en
saadan Optræden kunde vinde Tilslutning og betegne en ny Periodes Begyndelse?

Hvad var i det Hele Grunden til, at Humanismens Gennembrud i Slutningen af
det 18.\ Aarhundrede ikke kom til at indlede en Udvikling, der fortsattes i
lige Linie indtil nu, men afløstes af sin skarpe og bevidste Modsætning?

Da en nærmere psykologisk-historisk Forklaring af det vældige Omslag, der kom
til at udfylde saa stor en Del af vort Aarhundrede, og gjorde det til
»Modsætningernes Tid«, vil kunne tjene til nærmere Belysning af vort Problems
Natur, ville vi gaa noget nærmere ind derpaa, men saaledes, at vi især holde
os til de Aarsager, der kunne findes i selve Beskaffenheden og Ejendommeligheden
af det 18.\ Aarhundredes Humanisme. Der virkede naturligvis mange forskellige
Aarsager med. Saaledes fik den almindelige politiske Reaktion ogsaa Indflydelse
paa den religiøse Udvikling; Avtoritetsprincipets Genopstandelse i Staten maatte
ogsaa føre til dets Begunstigelse paa det religiøse Omraade, og de politiske
Bagstrævere benyttede sig ofte af Religionen paa den modbydeligste Maade. Men
dersom der ikke havde været dybt liggende indre Aarsager, vilde denne ydre
Indflydelse ikke have kunnet udrette saa Meget. ---

\medskip
\noindent 1. Enhver Reaktion skyldes for en stor Del en Kontrastvirkning.
Ligesom det af Lyset trættede Øje finder Hvile ved Mørket eller dog ved et
mattere Skær, saaledes kan i det Hele den menneskelige Aand ikke udholde en
stærk og langvarig Anspændelse af en enkelt Side i dens Natur, men kaster sig
uvilkaarligt over til den modsatte Side for at finde en ny Hvilestilling. Med
Hensyn til det Problem, vi her tale om, er det især Modsætningen mellem den
aktive og den passive, den virkende og den modtagende Side i vor Natur, som
bliver af Betydning. Ikke som om der i vort aandelige Liv kan paavises
Tilstande eller Fænomener, hvor vi enten forholde os absolut virkende eller
absolut modtagende: ved enhver Form for aandeligt Liv og i enhver aandelig
Tilstand kan der tvertimod paavises baade en Aktivitet og en Passivitet. Men
disse to Momenter kunne staa i et meget forskelligt indbyrdes Forhold; snart
har det ene, snart det andet Overvægten. Saaledes have den forskende og
grublende Tænken, den prøvende og fremadstræbende Samvittighed og den formende
og skabende Fantasi en overvejende aktiv Karakter, og i Modsætning til dem staa
i denne Henseende Trangen til Trøst og Hengivelse, og Følelser som Pietet og
Resignation. Hvor det gælder om at frembringe nye Former for Livet, faar den
aktive Side i vor Natur Overvægten og lægger Beslag paa hele Energien. En
saadan Opgave stillede sig netop for Slægten i den sidste Del af det
18.\ Aarhundrede. Det var den Tid, hvor ubegrænset Fremadskriden var Slagordet,
og hvor Træghed betragtedes som Arvesynden. Men der maatte komme en Tid, da
Menneskenaturens andre Sider krævede deres Tilfredsstillelse, og hvor man søgte
Hvile efter det rastløse Arbejde. Og Hvilen fandtes ikke i de nye Former. De
behøvede endnu alt for meget aktiv Medvirken, stillede stadigt Fordring om
Selvvirksomhed. Hvile kunde derimod maaske findes i de gamle Former, som man i
sin Fremstormen havde vendt Ryggen, men om hvilke man nu opdagede, hvor
vidunderligt de vare tilpassede efter Følelseslivets Trang: i Aarhundreder eller
Aartusinder havde Slægt efter Slægt nedlagt sine Erfaringer i dem, havde lagt
dem til Rette, saa der ikke behøvedes uafbrudt Arbejde for at bruge dem. Vor
Aktivitet søger det Nye; med vor virkende Natur leve vi i Nutiden; men den
passive og modtagende Side i os bestemmes og næres overvejende ved den
forudgaaende Tids Eftervirkninger; her nærmer det aandelige Liv sig det Ubevidste
og drager Næring af det. Et Eksempel herpaa fra det enkelte Individs Liv have vi
i Barndomserindringernes Magt: de virke uden vor bevidste Medvirkning og stige
især op med overvældende Styrke, naar Sindet er træt af aktivt Aandsarbejde.

Det er karakteristisk, at Schleiermacher netop ved Overgangen til det nye
Aarhundrede opstillede sin berømte Definition af Religionen som den ubetingede
Afhængighedsfølelse. Det var som et Program for den Retning, Aandslivet nu skulde
tage. Humanismen blev afløst af Konfessionalismen, den frie Forsken og den frie
Samvittighed af Avtoritetstroen.

Ere vi da henviste til et Sisyfusarbejde, ruller Stenen stedse tilbage, naar vi
vel have faaet den op paa Bjærget? --- Nej, men i det Anførte ligger, at kun den
Sandhed kan blive sejrende, der kan gennemtrænge og fylde hele Livet. Og det sker
ikke med ét Ryk; det tager Tid, og først gennem mange Svingninger skrider det
fremad. Det synes at være en psykologisk Lov, at Erkendelsen udvikler sig
hurtigere end Følelseslivet. Kun langsomt nedfældes det aktive og bevidste
Arbejdes Resultater i Følelsens uvilkaarlige Liv. Før det kan ske, kommer der da
naturligt en Periode, i hvilken de nye Tanker blegne og vise sig utilfredsstillende;
Tvivl og Uro bemægtiger sig Mange, og der synes ikke at være nogen anden Udvej end
at kaste sig tilbage i Følelseslivets gamle Former, hvis man ikke skal gaa til
Grunde. Men denne Disharmoni vil kun være forbigaaende: ere de nye Tanker gyldige,
udtrykke de virkeligt nogle af de Love, vor Tilværelse er underlagt, saa ville de
stedse arbejde sig frem paa ny, og den ustadige Ligevægts Tid vil være forbi. I
Længden vil Følelseslivet udvikle sig i større eller mindre Harmoni med Tankelivet:
paa den Maade bleve jo selve hine gamle Former og Dogmer til. Allerede selve
Reaktionsperioderne vise sig at bære Præg af det, de bekæmpe; i de gamle
Læderflasker vil der for en ikke ringe Del hældes ny Vin, og kun den, der holder
sig til det Udvortes, vil mene, at Alt igen er blevet som det var før.

Vi staa her over for en aandelig Natursammenhæng, som vi maa bøje os for, og som
vi ikke kunne trodse. Den, der har faaet ret Øje for den, vil hverken have Manges
let købte Tillid til Sejrens Nærhed eller tabe Modet ved tilsyneladende Tilbagegang.
Han véd, at ligesom de tilsyneladende pludselige Revolutioner kun betegne, at det
længe i Skjul Forberedte nu træder frem for Dagens Lys, saaledes tyde de pludselige
Reaktioner paa, at der var Drifter og Fornødenheder, som kun ufuldkomment bleve
tilfredsstillede ved det tilsyneladende saa glimrende Fremskridt, og som derfor maa
søge Tilfredsstillelse ved det Gamle, indtil de have lempet sig efter det Nye, eller
dette efter dem.

\medskip
\noindent 2. Det 18.\ Aarhundredes Humanisme gav ikke blot den erkendende og
virkende Del af vor Natur en afgjort Overvægt over den følende og modtagende, men i
nøje Sammenhæng hermed havde den en afgjort aandsaristokratisk Karakter. Den hele
Bevægelse var, og maatte efter sin Natur være, indskrænket til en lille Kreds; det
egentlige Folk paavirkedes saare lidet af den. Til ingen Tid have maaske de vel
stillede Klasser i Samfundet vist saa stor Iver for at gaa i Spidsen for Oplysningens
og Fremskridtets Arbejde. Men Afstanden mellem de Ledende og dem, der trængte til
Ledelse, var overordentlig stor. Og Udviklingen havde ført til et Brud med den fælles
Tro, de fælles Former for de højeste Tanker. I Middelalderen kunde den dybest
grublende Skolastiker og den enfoldigste Træl mødes i samme Bekendelse. Nu var denne
Enhed sprængt, og der maatte, især da samtidigt Følelsen for det Nationale og
Folkelige tog et stærkt Opsving, opstaa en Trang til ogsaa i Tro at være Et med sit
Folk. Reaktionen mod det 18.\ Aarhundrede fik fra denne Side set en aandsdemokratisk
Karakter. Det var jo --- for at tage et enkelt Eksempel --- den Overbevisning, at
»de Ulærdes Tro umuligt kan være bunden til de Lærdes Vidnesbyrd«, der førte Grundtvig
til hans »Opdagelse«.

Ligesom Spliden mellem de forskellige Sider eller Kræfter i det aandelige Liv vil
ogsaa Spliden inden for Folkets Liv være et stadigt virkende Motiv til at vende
tilbage til den gamle Tro. Fælles Former og Udtryk for de højeste Tanker skabes
langsomt og forudsætte, at Tankerne ere gaaede over i Kød og Blod hos Alle. Da nu
de individuelle Forskelligheder i religiøs Trang og religiøs Sans maa antages at
stige og tiltage, jo mere Religionsfriheden bliver til Virkelighed, saa vil et saadant
Fællesskab blive stedse vanskeligere at tilvejebringe.

Et Vidnesbyrd om, hvor stor Trangen til Fællesskab kan være, frembyder Schellings og
Hegels Religionsfilosofi. Disse Mænd og deres Tilhængere vare overbeviste om, at de i
deres filosofiske Ideer havde udviklet det samme Indhold, som Kirkens Dogmer udtalte.
Forskellen skulde kun angaa Formen: hvad Teologien har som Forestilling, i Fantasiens
og Følelsens Form, det skulde Filosofien forme til abstrakt Begreb, uden at derved det
Væsentlige gik tabt. I Dogmet om Guds Menneskeblivelse f.\ Eks.\ skulde indeholdes den
filosofiske Tanke, at selv det Højeste maa udvikle sig gennem Kamp og Modsætninger, og
at et Princip kun viser sin Gyldighed, naar det formaar at gennemtrænge det virkelige
Liv med al dets Splid og med alle dets stridende Forhold.

I og for sig er det fuldt berettiget at betragte Dogmerne som Symboler. For den
Ikke-troende kunne de ikke være Andet, naar de i det Hele skulle være Noget. Og Mange,
der ansé sig selv for meget troende og betragtes af Andre som saadanne, staa i
Virkeligheden paa dette Standpunkt. Men det vil altid føre til mislige Følger, naar man
benytter den symbolske Opfattelse til at tilvejebringe et falskt Skin af Fællesskab. En
virkelig Enhed kan ad den Vej ikke bringes til Veje. Paa et saadant falskt Skin beroede
det sidste Forsøg paa Forsoning af Tro og Viden, som er blevet gjort, Hegels
Religionsfilosofi, af hvilken Mange ventede sig en »ny Æra«. Vergil skulde her ikke ydmyg
underordne sig Beatrice; de skulde følges Arm i Arm. Men denne tilsyneladende Enighed gav
kun Anledning til, at Modsætningen mellem Tro og Viden blev udtalt i endnu langt skarpere
og mere tilspidset Form end hidtil, saaledes som det skete af Strauss og Ludwig Feuerbach i
Tyskland og af Søren Kierkegaard hos os. Disse Mænd vare enige i, at den teologiske Tros
Indhold var det Ubegribelige; Forskellen mellem dem bestod i, at Strauss og Feuerbach sagde:
Det er ubegribeligt, og derfor kan det ikke tros! medens Kierkegaard sagde: Netop fordi det
er ubegribeligt, kan og skal det tros! --- Strauss's Ord: »Af falske Foreningsforsøg er der
sket nok; kun Skærpelse af Modsætningerne kan føre videre!« er karakteristisk for
Udviklingens Gang i det 19.\ Aarhundrede. En berømt kirkehistorisk Forsker, F.\ C.\ Baur,
har med Rette fundet vort Aarhundredes Ejendommelighed deri, at alle Modsætninger --- i
Videnskab, Religion og Politik --- udarbejdes til den højeste mulige Grad. Det hænger
naturligt sammen med, at Reaktion og Revolution paa de forskellige Omraader saa hyppigt
have staaet over for hinanden; derved er der opstaaet en klar Bevidsthed om Modsætningernes
Natur og en Bestræbelse for at udforme dem hver for sig i deres fulde Konsekvens. Baade
Nødvendigheden af at træffe et Valg mellem uforenelige Modsætninger og Karakteren af de
Veje, man skal vælge mellem, gaar mere og mere op for den almindelige Bevidsthed. ---

Har den blotte Trang til Fællesskab saaledes ført til Sympati med den gamle Tro --- hvad
enten denne Sympati førte til personlig Tilslutning til dens Bogstav eller blot til
symbolsk Udlægning af dens Udsagn, --- saa maatte dette hele Spørgsmaal om Tro og Viden
særligt stille sig paa en alvorlig Maade for dem, der indsaa dets Sammenhæng med det sociale
Spørgsmaal. Selv om man indrømmer, at den videnskabelige Kritik af Dogmerne fører til deres
Opløsning, vil dog den Tanke rejse sig, at disse Dogmer, saa fornuftstridige de end maatte
være, ere Manges eneste aandelige Tilhold og Trøst, --- det, der holder deres Haab oppe under
Livets tunge Kamp, --- det Eneste, de raade over, der kan løfte deres Blik ud over den trange
Kreds, i hvilken deres Liv bevæger sig. Fra dette Synspunkt har saaledes en af vor Tids
ædleste filosofiske Tænkere, Albert Lange, Forfatteren til »Geschichte des Materialismus«,
udtalt sig over for Strauss og Feuerbach, med hvem han er enig hvad den rent teoretiske Side
af Sagen angaar. Det er i det Hele karakteristisk, hvor stor en Rolle Spørgsmaalet om
Religionens Nytte og Uundværlighed spiller i den religiøse Debat. Det gælder i første Linie
ikke: Sandhed eller ikke Sandhed? --- men: Trøst eller ikke Trøst? --- I tidligere Tiders
Debat kom denne Side af Sagen ikke frem som Noget for sig, end sige som Hovedsagen. Man søgte
ikke efter Grunde og Motiver til at holde paa den dogmatiske Tro trods Vanskelighederne; men
det var dem, der opgav den dogmatiske Tro, hvem det paalaa at motivere deres Skridt.

Det er sikkert af meget stor Betydning, at det aandelige Fællesskab ikke forsvinder, og at
vor Slægt ikke mister nogen Kraft, der kan holde den oppe i Kampen for Tilværelsen. Det er
med fuld Ret, at disse Hensyn blive bestemmende for den Maade, paa hvilken den religiøse Kamp
føres; men selve Kampen kunne de kun for en Tid dæmpe, da de ikke anvise nogen Vej til
Problemets Løsning.

\medskip
\noindent 3. Endelig havde den humane Livsanskuelse, saaledes som det 18.\ Aarhundrede
formaaede at fremstille den, en ensidig abstrakt og idealistisk Karakter. Den udtaltes efter
sit almindelige Princip og med de vigtigste Forudsætninger og Fordringer, den medfører, men
savnede tilstrækkelig Gennemførelse og realt Grundlag. Saa at sige ethvert Fremskridt, det
19.\ Aarhundrede har gjort, har bidraget Sit til at raade nogen Bod paa disse Mangler.
Naturvidenskabernes storartede Opsving har givet en fastere Basis for Overbevisningen om den
lovordnede Sammenhæng i Naturen. I Udviklingens Begreb er der paavist en ledende Idé, om
hvilken al vor Viden paa Naturens og Historiens Omraade kan samle sig. Den historiske Kritik
og den sammenlignende Religionsvidenskab have udvidet og berigtiget det Materiale, som er
nødvendigt til religionsfilosofiske Problemers Behandling. En dybere gaaende psykologisk
Forstaaelse af det menneskelige Aandslivs forskellige Sider i deres indbyrdes Forhold og en
fornyet Drøftelse af de erkendelsesteoretiske Spørgsmaal giver Problemet om Tro og Viden en
noget anden Karakter end det forhen har kunnet have. --- Alt dette er naturligvis kun smaa
Fremskridt i Forhold til Problemets Storhed. Dette vil i lange Tider følge Slægten paa dens
Vej, og den forskellige Maade, hvorpaa det opfattes og behandles, vil kunne afgive Maalestok
for den aandelige Udviklings Grad og Retning til de forskellige Tider. Her have vi kun tilbage
at vise, hvorledes Problemet maa stille sig i vore Dage.

\section*{VI.}

Der blev i den Forhandling om Tro og Viden, som i Tredserne fandt Sted i vor
Literatur, lagt stor Vægt paa Distinktionen mellem Tro og Teologi, og jeg selv
hørte, som Discipel af Rasmus Nielsen, til deres Tal, der knyttede ikke ringe
Forhaabninger til den. Den hjælper os aabenbart ikke, og den er i hvert Tilfælde
ikke holdbar. Saa snart Troen har et bestemt Indhold, en bestemt Genstand, maa
der ogsaa udvikle sig en Lære om dette Indhold, som søger at paavise dets indre
Sammenhæng og dets Betydning. Ellers véd hverken den Troende selv eller Andre,
hvad det egentligt er, han tror paa.

Nu til Dags har desuden Teologien opgivet at støtte sine Paastande med
videnskabelige Beviser. Den har opgivet alle stolte Drømme om at faa Vergil til
at gaa i Ledebaand og er tilfreds, naar man lader den uanfægtet fra den
videnskabelige Kritiks Side. Den erklærer sine Dogmer for ubevislige, indrømmer
maaske endog, at de frembyde Vanskeligheder for den »naturlige« Forstand, og
stiller sig væsentligt paa det Standpunkt, at den hævder den dogmatiske Tro som
en aandelig Livsfornødenhed --- vel at mærke en Fornødenhed, det er Pligt at
have. Det skal være Tegn paa Forhærdelse og Trods mod det Hellige, naar man ej
føler den Trang, der lader Fornuften bøje sig for det Ubegribelige. Deri ligger,
som jeg allerede har haft Lejlighed til at berøre, Problemets egentlige Brod.

Paa den anden Side vil ingen videnskabelig Tænker paastaa, at den absolute
Umulighed af Dogmernes Indhold kan strengt bevises. Vi bygge i videnskabelig
Forsken paa de naturlige Aarsagers Princip. Det har hjulpet os igennem mangen god
Gang, og noget andet videnskabeligt Princip til Forstaaelse af Tilværelsen vil
neppe nogen Sinde kunne opstilles; --- men det kan ikke bevises, at Alt i
Tilværelsen er dette Princip underlagt.

Ja, vi maa endog gaa et Skridt videre. Samtidigt med de vidunderlige Fremskridt,
som ere gjorte i Kundskab om Naturen og Historien, baade hvad Enkeltheder og hvad
den store Sammenhæng angaar, have vi i langt dybere Forstand end forhen lært, at
den hele Tilværelse er én stor Gaade. For Filosofen er det ringeste Fænomen i
denne Henseende lige saa lærerigt som det største og mest i Øjne faldende. Det
simpleste Aarsagsforhold giver os Mere end nok at tænke paa --- indeslutter
egentligt hele Verdensgaaden i sig. Alt hvad vi forstaa og forklare, forstaa og
forklare vi ved at paavise Aarsagssammenhæng: vidste vi nu blot, hvad
Aarsagssammenhæng egentligt er, --- hvad det er, der holder Alt i Naturen sammen
og gør den til den vidunderlige Helhed, der mere og mere afslører sig for os, ---
saa vilde alle Gaader være løste! Men til en saadan Kundskab er al Udsigt spærret.
Vore Fornemmelser, Forestillinger og Begreber (og saaledes ogsaa Aarsagsbegrebet)
ere ret gode Hjælpemidler til at finde os til Rette i Tilværelsen, men vi have
ingen Ret til at mene, at de kunne afbilde Tilværelsens inderste Væsen for os. Vi
kunne i det Højeste forme spekulative Hypoteser om, hvorledes Tilværelsen i sig
selv kunde være beskaffen; men nogen Erfaringsbekræftelse kunne saadanne Hypoteser
i Følge Sagens Natur ikke faa.

Deraf følger nu aldeles ikke, at man bøjer sig for ethvert teologisk Orakelsprog.
Deraf følger kun, at Diskussionen om Tro og Viden maa føres paa en noget anden
Maade end den hidtil er bleven ført. Lad os i Korthed berøre de Hovedspørgsmaal,
som her ville være at drøfte.

\medskip
\noindent 1. Naar den teologiske Tro vedkender sig, at dens Paastande ere
ubevislige eller maaske endog uforenelige med nogle af Videnskabens vigtigste
Forudsætninger og Resultater, stiller den sig tilsyneladende paa et utilgængeligt
og uangribeligt Standpunkt. Der synes ikke mere at være nogen fælles Grund, paa
hvilken Kampen kan føres.

Men naar den teologiske Tro alligevel --- for at udtrykke sit Indhold --- benytter
Forestillinger og Begreber, som ogsaa den naturlige Fornuft og den videnskabelige
Tænkning anvender, saa maa der spørges, om disse Begreber tages i aldeles samme
Betydning. Det har vist sig forbundet med uovervindelige Vanskeligheder, selv for
den dristigste Spekulation, at gennemføre Sætninger og Domme om det, der ligger ud
over enhver mulig Erfaring. Hvad kan det nu egentligt være, som bevirker, at Troen
her kommer bedre fra det end den videnskabelige Tænken? Hvorfor bortfalde de store
Vanskeligheder ved Dogmerne, naar man gaar fra Viden over til Troen? En Skabelse af
Intet f.\ Eks., --- en første Aarsag, --- et absolut godt Væsen, som dog er Aarsag
til Alt, altsaa ogsaa til det Onde, --- saadanne Begreber, mod hvilke den menneskelige
Tanke er tørnet atter og atter, indtil den til sidst har opgivet Forsøgene paa at
forene dem med sine andre Begreber, --- hvorledes kunne de stille sig i et helt andet
Lys, fordi de erklæres at være Genstande for Tro, ikke for Viden?

Svares der (som f.\ Eks.\ af Biskop Martensen i hans Dogmatik), at den Genfødelse
af Menneskenaturen, som finder Sted hos den Troende, ogsaa omfatter Erkendelsen,
saa at det bliver forstaaeligt, der for var ubegribeligt eller endog selvmodsigende,
saa kunne naturligvis de Ikkegenfødte ikke tale med om, hvorledes Tingene tage sig
ud efter en saa radikal Forandring. Men da de Troende og Genfødte dog ikke ganske
opgive Sammenhængen med den naturlige Verden og, for saa vidt de vedblive at høre
med til den, ogsaa maa vedblive at bruge deres naturlige Fornuft, saa maatte de dog
kunne sammenligne de to Arter af Erkendelse, den naturlige og den genfødte, og sige
os, hvori Forskellen egentligt bestaar. Her er der et Kapitel i Dogmatiken, som
Teologerne i alle de mange Aar, i hvilke Teologien har floreret, have forsømt at
skrive, nemlig: en ny eller højere Logik og Erkendelsesteori, som kunde vise os, ved
Hjælp af hvilke nye Principer man i Troen slipper forbi de Skær, paa hvilke man i
den naturlige Tænkning bestandigt strander. Indtil dette Kapitel er skrevet, maa vi
have Ret til at undersøge de teologiske Begreber paa samme Maade, som vi undersøge
alle andre, og fælde samme Dom om Inkonsekvenser og uberettigede Begrebsanvendelser,
som vi alle andre Steder fælde. En saadan Undersøgelse vil føre til det Resultat, at
Dogmet i Virkeligheden staar lige saa afmægtigt over for de store Gaader som
Videnskaben.

\medskip
\noindent 2. Men --- vil man da sige fra teologisk Side (især nu, da den filosofiske
Interesse i vor teologiske Verden synes at være skrinlagt med Biskop Martensen) ---
det er dog en Misforstaaelse at mene, at Troen skal fyldestgøre en teoretisk Trang.
For Troen har Tænkning og Refleksion Intet at betyde, og al Verdens Modgrunde falde
magtesløse til Jorden af den simple Grund, at Ingen kommer til Troen uden den, som
søger Trøst for en ængstet Samvittighed. Det er Følelsen, særlig den ængstede
Samvittighed, ikke nogen anden Kraft, der her virker.

Dette er ikke nogen behagelig Indvending, hvis Meningen er, at den Talende selv ikke
kan undvære den Trøst, der ligger i Troen; thi i saa Fald maa Diskussionen høre op.
Hvem vil berøve Nogen hans Trøst? Selv den ivrigste Kemiker vilde ikke bruge til
Undersøgelse den Drik Vand, hvormed et Menneske skal læske sin brændende Tørst.
Enhver Diskussion forudsætter, at Deltagerne hver for sig have saa megen Ro og Kraft
i Sindet, at de kunne se Sandheden under Øjne, saa man ikke behøver, som over for
Syge, at dølge Noget.

Hvad der trøster den Ene, trøster desuden ikke den Anden, eller vækker endog hans
Forargelse og Indignation. Saa snart der appelleres til Følelsen, appelleres der mere
eller mindre til noget rent Individuelt. En saadan Appel kan aldrig være Argument;
thi Følelsen i og for sig lærer os aldeles Intet uden vor egen Tilstand. Den mest
levende Glæde og Beroligelse, jeg føler ved en Efterretning, garanterer Intet om
denne Efterretnings Vederhæftighed. Fra min Følelse kan jeg derfor ikke med Sikkerhed
slutte til Noget, som ligger ud over min egen Tilstand. Vore Ønsker bestemme ikke,
hvad der er Sandhed; men det kunde jo være, at H.\ C.\ Ørsted havde Ret i, at den
fulde Sandhed altid bringer sin Trøst med sig.

Den Følelse, paa hvilken man især beraaber sig som Støtte for den teologiske Tro, er
Syndsbevidstheden. Her opstaar nu det store Spørgsmaal, om Syndsbevidstheden,
saaledes som Teologien opfatter den, virkeligt er en, etisk set, sund og berettiget
Følelse, --- om der ikke skulde indeholdes Elementer i den, som ikke ere virkeligt
etiske, men som hænge sammen med barbariske Forestillinger (om en vred Guddom, der
skal forsones), og som maa udskilles, for at Syndsbevidstheden kan blive til virkelig
etisk Angerfølelse.

Men det er netop Maalestokken for det Etiske, der er saa forskellig paa de to
Standpunkter! Og Teologien, der før viste os fra det Dogmatiske over til det Etiske,
hævder nu omvendt, at uden Dogmetro er ingen Etik mulig. Appellen til
Syndsbevidstheden leder os derfor kun i en Kreds: thi Syndsbevidstheden er en af
bestemte dogmatiske Forestillinger betinget Følelse, og det er da intet Under, om den
aflægger Vidnesbyrd til Fordel for Dogmet. --- Vi komme kun ud af denne Kreds, dersom
Teologien kan vise os, hvad der hidtil ikke er lykkedes, at kun dogmatisk Tro kan
være Grundlag for Etik, ja, at de teologiske Antagelser i det Hele taget egne sig til
at være Grundlaget for en Etik. Dersom derimod en human Etik er mulig, og dersom en
ubevidst human Etik skulde vise sig stedse at ligge til Grund for den teologiske Etik
i dens forskellige Former (saaledes som man kunde slutte af den Kendsgerning, at naar
Menneskene forbedre sig, forbedre ogsaa deres Guder sig), saa vilde den skarpeste Brod
være taget bort, og hvad der saa blev tilbage af Problemet, vilde vel stedse kunne
sætte Følelsen i stærk og dyb Bevægelse, men ikke længer spalte vor Slægt i to
modsatte Lejre, der støde sammen i uforligelig Strid om den sidste Maalestok for Godt
og Ondt.

\medskip
\noindent 3. Der er en Paastand, som længe utaalmodigt har ventet paa at komme til
Orde fra teologisk Side. Den lyder omtrent saaledes: »Troen er bygget paa historiske
Kendsgerninger; dens Indhold er først og fremmest Historie og hverken noget udspekuleret
System eller blotte Følelsespostulater. Og historiske Kendsgerninger maa man tage, som
de ere; ved dem kan intet Fornuftræsonnement rokke.«

Naar man nu blot her tager Ordet Historie i samme Betydning, som vi tage det i den
historiske Videnskab! Men det gaar med den teologiske Brug af dette Begreb som med
Brugen af andre Begreber (Aarsag, Etik o.\ s.\ v.): trods Ordligheden er den virkelige
Betydning af Begreberne saare forskellig. Historisk er ikke Alt hvad der staar i gamle
Bøger eller indeholdes i gamle Traditioner. En historisk Kendsgerning kan ofte kræve
indtrængende kritisk Forsken, før den kan gøre Fordring paa Anerkendelse. Det store
Spørgsmaal bliver da, om Teologien vil underkaste sig den historiske Kritiks sædvanlige
Love, eller om den gør Fordring paa visse Privilegier for sine »Kendsgerninger«. At den
maa fordre saadanne Privilegier, følger allerede af, at dens Kildeskrifter indeholde
Fortællinger om Undere, der kræve Tro for at blive anerkendte. En historisk Forsken,
der ikke anvender andre Forudsætninger end dem, som al Videnskab maa gaa ud fra, fører
--- som Erfaring stedse har vist --- til en helt anden Opfattelse af Begivenhederne ved
Kirkens Grundlæggelse end den, som den teologiske Tro fastslaar. Det gaar altsaa med
Appellen til Historien som med Appellen til Følelsen og til Etiken: den fører i en
Kreds, thi ved Historie mener man den teologisk opfattede og vurderede Historie. ---

\end{document}
%%% Local Variables:
%%% mode: latex
%%% TeX-master: t
%%% End:
